\section{Introduction}

Smoothing parameters are often introduced as technical constants, although in forecasting they determine a substantive modeling choice: how much recent movement is treated as signal and how much is treated as noise. A small penalty lets the fitted trend track the observations closely, while a large penalty forces the fitted path toward a low-degree polynomial. If that trend is subsequently extrapolated, the smoothing choice also changes the future trajectory.

Finite-difference penalized least squares (PLS) makes this trade-off especially transparent. Guerrero \cite{guerrero2007penalized,guerrero2008smoothness} showed that the penalty can be translated into an interpretable percentage of smoothness. This controlled-smoothness view is useful because the same numerical penalty does not represent the same effective amount of smoothing across sample sizes and difference orders. The remaining practical question is immediate: if the fitted trend will be used for forecasting, which percentage of smoothness should be chosen?

The forecasting literature provides strong precedents for choosing tuning parameters by predictive performance. Forecast-error minimization has long been used to estimate exponential-smoothing parameters \cite{taylor2004stes}; time-series cross-validation selects a kernel bandwidth through one-step-ahead prediction in dependent data \cite{hart1994tscv}; and recent work in the \emph{Journal of Forecasting} continues to connect hyperparameter selection, trend filtering, and out-of-sample loss to forecast performance \cite{zafar2022filter,stanek2023oos,wolff2024stock,franjic2025predictor,xu2025vix}. These papers motivate forecast-oriented tuning, but they do not define the smoothness percentage of a finite-difference penalized trend through a horizon-specific rolling future-block loss.

For fixed difference order $d$, estimation-window length $L$, and forecast horizon $h$, let $S\in[0,1]$ denote normalized smoothness. At each historical forecast origin, the trend is fitted using only observations available at that origin, extrapolated by the continuation implied by the same finite-difference model, and scored against the next $h$ realized observations. Averaging these errors yields a chronological criterion $F_{d,L,h}(S)$. We define
\begin{equation}
\label{eq:main_target_intro}
S^\star_{d,L,h}
\in
\arg\min_{S\in[0,1]}F_{d,L,h}(S)
\end{equation}
as the \emph{forecast-optimal smoothness} for that forecasting configuration.

The proposal is deliberately narrower than a general automatic forecasting system. The order $d$, window length $L$, and horizon $h$ are treated as fixed when defining the continuous smoothness problem. Joint adaptation of all three quantities is a different model-selection problem. Likewise, the present paper defines and studies the forecast target; a companion numerical study treats the separate problem of efficiently recovering multiple minima when the one-dimensional objective is non-unimodal.

The distinction from existing smoothing criteria is also important. Classical cross-validation and generalized cross-validation are usually motivated by estimation or reconstruction performance, while information criteria optimize different statistical objectives. Cort\'es-Toto et al.\ \cite{cortes2017smoothness} study the smoothness achieved by several such selectors inside PLS. The criterion proposed here is not intended to dominate those methods universally. It answers a different question: conditional on using the penalized trend as a forecasting device with a specified continuation and horizon, which amount of smoothness minimizes chronological forecast loss?

The contribution has four parts. First, controlled smoothness is converted into an endogenous forecast-based selection rule. Second, the normalized $S\in[0,1]$ coordinate makes the unsmoothed and maximally smoothed finite-difference limits explicit. Third, the target is horizon-specific through genuine future blocks rather than only one-step reconstruction or leave-one-out error. Fourth, tuning is embedded in a strictly chronological rolling-origin protocol, so the information set matches actual forecast use.

The remainder of the paper positions the proposal relative to controlled smoothness, predictive smoothing selection, and forecast evaluation; develops the finite-difference PLS representation; defines the forecast-optimal criterion; establishes its basic properties; and specifies the simulation and empirical protocol required to evaluate when forecast-based smoothness differs materially from conventional alternatives.
